\documentclass[11pt,a4paper]{article}

% === 中文支持 ===
\usepackage{ctex}
\usepackage[T1]{fontenc}

% === 页面布局 ===
\usepackage[top=2.5cm, bottom=2.5cm, left=2.5cm, right=2.5cm]{geometry}

% === 表格 ===
\usepackage{booktabs}
\usepackage{longtable}
\usepackage{array}
\usepackage{enumitem}

% === 颜色 ===
\usepackage{xcolor}
\definecolor{primary}{HTML}{404040}
\definecolor{success}{HTML}{22c55e}
\definecolor{danger}{HTML}{e74c3c}
\definecolor{warning}{HTML}{f59e0b}

% === 数学符号 ===
\usepackage{amssymb}

% === 超链接 ===
\usepackage[hidelinks]{hyperref}
\hypersetup{colorlinks=true, linkcolor=primary, urlcolor=blue!60!black}

% === 标题样式 ===
\usepackage{titlesec}
\titleformat{\section}{\Large\bfseries\color{primary}}{}{0em}{}[\titlerule]
\titleformat{\subsection}{\large\bfseries\color{primary}}{}{0em}{}

% === 文档信息 ===
\title{\textbf{DEMAND-001：角色分工重新定义}\\[0.5em]\small 需求变更文档 v0.1}
\author{万能产品小助手（PM AI）\\魏子政 审阅}
\date{2026-05-06}

\begin{document}

\maketitle

\begin{center}
\begin{tabular}{ll}
\toprule
\textbf{提出日期} & 2026-05-06 \\
\textbf{提出者} & 魏子政 \\
\textbf{优先级} & \textcolor{danger}{P0} \\
\textbf{状态} & \textcolor{success}{已完成} \\
\textbf{影响范围} & AGENT-ARCHITECTURE.md, TEAM-MANUAL.md, skills/pm/SKILL.md \\
\bottomrule
\end{tabular}
\end{center}

\vspace{1em}

\section{需求描述}

魏子政明确重新定义 OpenClaw 和 Cursor 的分工：

\begin{quote}
``你主要负责需求生成拆解和验收。Cursor 负责前后端（包含 UI 设计，因为 UI 也包含在前端里）、环境、运维、测试（和验收不一样！Cursor 需要根据自己生成的内容自验收，也需要回应你验收的需求）。''
\end{quote}

核心变更点：
\begin{enumerate}
  \item PM（OpenClaw）职责收窄：只做需求生成、需求拆解、独立验收
  \item Cursor 职责扩大：前端（含 UI 设计）+ 后端 + 环境 + 运维 + 测试
  \item 测试不等于验收：Cursor 做自验收 + 响应 PM 验收；PM 做独立验收
  \item UI 设计归 Cursor：不再由 PM 设计方案再交给 Cursor 实现
  \item 需求变更文档：每次变更后生成的技术文档放到 demands/ 目录，使用 tex-pdf 格式
\end{enumerate}

\section{背景与动机}

之前的分工（v1.0）中，PM 同时承担了 UI 设计方案输出、功能测试等职责，导致：
\begin{itemize}
  \item PM 承担过多执行层面的工作，偏离产品经理核心职责
  \item UI 设计由 PM 出方案、Cursor 实现的模式效率低，Cursor 本身有能力自主设计
  \item 测试和验收职责混淆，缺乏独立的验收视角
\end{itemize}

重新定义后，PM 聚焦需求侧（生成需求、拆解任务、独立验收），Cursor 聚焦实现侧（设计 UI、写代码、搭环境、自测、响应验收），两者职责清晰，减少协作摩擦。

\section{详细方案}

\subsection{新分工矩阵}

\begin{table}[h]
\centering
\begin{tabular}{lcc}
\toprule
\textbf{职责} & \textbf{OpenClaw (PM)} & \textbf{Cursor Agent} \\
\midrule
需求生成   & \textcolor{success}{\checkmark} & \\
需求拆解   & \textcolor{success}{\checkmark} & \\
UI 设计    & & \textcolor{success}{\checkmark} \\
前端开发   & & \textcolor{success}{\checkmark} \\
后端开发   & & \textcolor{success}{\checkmark} \\
环境搭建   & & \textcolor{success}{\checkmark} \\
运维部署   & & \textcolor{success}{\checkmark} \\
自验收     & & \textcolor{success}{\checkmark} \\
响应PM验收 & & \textcolor{success}{\checkmark} \\
独立验收   & \textcolor{success}{\checkmark} & \\
文档维护   & \textcolor{success}{\checkmark} & \\
进度管理   & \textcolor{success}{\checkmark} & \\
汇报沟通   & \textcolor{success}{\checkmark} & \\
\bottomrule
\end{tabular}
\caption{角色分工矩阵}
\end{table}

\subsection{自验收 vs 独立验收}

\subsubsection{Cursor 自验收}（每次开发完成后执行）：
\begin{itemize}
  \item 运行 \texttt{npm run build} 确认前端无报错
  \item 运行后端启动确认服务正常
  \item 对照 testing.mdc 逐项自检
  \item 如有数据库变更，确认迁移文件已生成
\end{itemize}

\subsubsection{PM 独立验收}（Cursor 返回结果后执行）：
\begin{itemize}
  \item 检查产出物是否完整覆盖需求
  \item 站在用户角度验证功能是否可用
  \item 执行查证三问（目标达成？细节硬伤？换位质疑？）
  \item 不依赖 Cursor 的自测结论，独立判断
\end{itemize}

\subsubsection{验收不通过的处理}：
\begin{enumerate}
  \item PM 记录问题到 BUGS.md
  \item 将问题反馈给 Cursor
  \item Cursor 修复后重新自验收
  \item PM 重新独立验收
  \item 循环直到通过
\end{enumerate}

\section{影响评估}

\subsection{文档影响}

\begin{table}[h]
\centering
\begin{tabular}{lp{8cm}}
\toprule
\textbf{文档} & \textbf{变更内容} \\
\midrule
AGENT-ARCHITECTURE.md & 分工矩阵重写，新增第六章 UI 辅助方案调研 \\
TEAM-MANUAL.md & 角色分工表重写，删除独立 UI 设计师角色 \\
skills/pm/SKILL.md & 职责描述收窄为需求生成+拆解+验收 \\
demands/README.md & 新建，定义需求变更文档规范 \\
CHANGELOG.md & 记录本次变更 \\
\bottomrule
\end{tabular}
\end{table}

\section{验收标准}

\begin{itemize}
  \item[$\square$] AGENT-ARCHITECTURE.md 分工矩阵中，UI 设计、环境、运维、自验收均标注为 Cursor
  \item[$\square$] TEAM-MANUAL.md 角色表删除独立的 UI 设计师行
  \item[$\square$] TEAM-MANUAL.md 明确区分自验收和独立验收
  \item[$\square$] demands/ 目录已创建，包含 README.md 文档规范
  \item[$\square$] CHANGELOG.md 记录本次变更
  \item[$\square$] 所有受影响文档版本号已更新
\end{itemize}

\vspace{1em}
\textit{以上验收项已于 2026-05-06 全部通过。}

\end{document}
